\section{Вимоги до логування та аудиту}\label{sec:loguvannia}

Система веде два незалежні рівні журналювання: (1) операційні логи бізнес-подій у процесі функціонування застосунку та (2) протоколи виконання автоматизованих тестів у процесі TDD~\sked{K8}. Усі вимоги мають ідентифікатори \texttt{REQ-LOG-xxx}.

\subsection{Операційне логування подій}\label{sec:loguvannia-podii}

\begin{itemize}[leftmargin=*,labelindent=0pt]
    \item \textbf{REQ-LOG-001 (Обов'язкові типи подій для логування):} Система повинна автоматично фіксувати у журналі такі типи подій:
    \begin{enumerate}
        \item Створення нового бронювання користувачем;
        \item Скасування бронювання користувачем;
        \item Модифікація довідників адміністратором (створення, оновлення або скасування фільмів, кінозалів і сеансів);
        \item Відмови в операціях через конфлікти зайнятих місць або порушення правил бронювання;
        \item Непередбачені системні збої та внутрішні помилки застосунку.~\sked{K16}
    \end{enumerate}
    \textit{Умова верифікації:} Кожна з перелічених дій ініціює запис відповідного рядка у журнал аудиту.

    \item \textbf{REQ-LOG-002 (Структура та обов'язковий склад полів запису):} Кожен запис операційного журналу повинен бути структурованим (наприклад, JSON-об'єкт) та містити обов'язкові поля, наведені у \tabref{tab:log-fields}~\sked{K17}.
\end{itemize}

\begin{table}[htbp]
\caption{Обов'язкові атрибути операційного лог-запису}\label{tab:log-fields}
\centering
\begin{tabularx}{\textwidth}{@{}l l L@{}}
\toprule
\textbf{Поле} & \textbf{Тип даних} & \textbf{Опис та призначення} \\
\midrule
\texttt{timestamp} & ISO-8601 UTC & Точна дата і час фіксації події з мілісекундами. \\
\texttt{event\_type} & String (Enum) & Тип події (\texttt{BOOKING\_CREATED}, \texttt{SESSION\_CANCELLED} тощо). \\
\texttt{status} & String & Результат операції (\texttt{SUCCESS}, \texttt{REJECTED}, \texttt{ERROR}). \\
\texttt{user\_id} & Integer / UUID & Внутрішній ID користувача або адміністратора (якщо ідентифікований). \\
\texttt{entity\_ids} & Object & Ідентифікатори об'єктів (ID бронювання, ID сеансу, номери місць). \\
\texttt{reason} & String & Детальна текстова причина відмови для неуспішних операцій. \\
\bottomrule
\end{tabularx}
\end{table}

\subsection{Захист конфіденційності та правила фільтрації логів}\label{sec:loguvannia-bezpeka}

\begin{itemize}[leftmargin=*,labelindent=0pt]
    \item \textbf{REQ-LOG-003 (Заборона збереження конфіденційних даних):} Система у версії 1 не збирає і не обробляє платіжні реквізити. До журналів логів суворо заборонено записувати паролі у відкритому чи хешованому вигляді, токени доступу (сесійні токени/JWT) та будь-які інші конфіденційні автентифікаційні дані~\sked{K18}.
    \textit{Умова верифікації:} Автоматизований тест-фільтр перевіряє вихідний потік логів під час реєстрації та входу; поля паролів замінюються маскою \texttt{[REDACTED]} або повністю вилучаються.
\end{itemize}

\subsection{Персистентність та зберігання логів для діагностики}\label{sec:loguvannia-zberezhennia}

\begin{itemize}[leftmargin=*,labelindent=0pt]
    \item \textbf{REQ-LOG-004 (Персистентність операційних логів):} Логи повинні записуватися на персистентний носій і зберігатися після штатного чи аварійного завершення роботи програми для можливості подальшого аналізу та діагностики помилок. Конкретний механізм збереження (структуровані файли, ротація, база даних) обирається під час проєктування технічної реалізації~\sked{K19}.
    \textit{Умова верифікації:} Після примусового перезавантаження серверного процесу журнал логів попередніх операцій зберігає цілісність і доступний для читання.
\end{itemize}

\subsection{Журналювання результатів автоматизованого тестування}\label{sec:loguvannia-testy}

\begin{itemize}[leftmargin=*,labelindent=0pt]
    \item \textbf{REQ-LOG-005 (Збереження протоколів запусків тестів TDD):} Результати кожного запуску автоматизованих тестів повинні зберігатися в окремих журналах (тест-логах). Журнал запусків фіксує: час початку й завершення тестування, загальну кількість тестів, повний перелік успішно пройдених і впалих тест-кейсів, а для кожного падіння --- детальний трасування стеку та діагностичне повідомлення про причину падіння~\sked{K20}.
    \textit{Умова верифікації:} Тестовий раннер автоматично генерує стандартизований звіт після кожного тестового прогону.
\end{itemize}
